\partframe{九}{补编：物理上怎么理解几何 Langlands}

\begin{frame}{一个自然的问题}
  Witten 的演讲到上一帧就结束了。但它留下一个没有正面回答的问题：
  \vspace{0.5em}
  \begin{center}
  \large 几何 Langlands 对应为什么应该是真的？
  \end{center}
  \vspace{0.5em}
  这不是“有没有证据”的问题，而是“它是不是某个更基本结构的影子”。
  2006 年，Kapustin 与 Witten 给出了一个答案；2008 年，Frenkel 与 Witten
  把它推进到内窥（endoscopy）。这一部分补上这条线。
  \begin{beizhu}
  \textbf{批注}\quad 以下不属于 Witten 2021 年演讲的原文，是译者根据
  Kapustin--Witten（2006）与 Frenkel--Witten（2008）补写的“下半场”。
  \end{beizhu}
\end{frame}

\begin{frame}{四维理论，二维约化}
  Kapustin--Witten 的出发点：取四维 $\cN=4$ 超杨--米尔斯理论，放在
  \[
    \RR\times\Sigma\times C
  \]
  上，其中 $C$ 是 Riemann 面，$\Sigma$ 是另一个曲面。先做拓扑扭化，再在 $\Sigma$ 上约化。
  \vspace{0.4em}

  结果是一个\textbf{二维 sigma 模型}，靶空间是 $C$ 上 $G$ 的 Hitchin 模空间：
  \[
    \mathcal M_H(G,C)=\Bigl\{(E,\phi)\;\Big|\;
    \bar\partial_A\phi=0,\; F_A+[\phi,\phi^*]=0\Bigr\}\Big/\!\sim .
  \]
  这里 $E\to C$ 是全纯 $G_\CC$-丛，$\phi\in H^0(C,\ad E\otimes K_C)$ 是 Higgs 场。
\end{frame}

\begin{frame}{Hitchin 模空间：三个复结构}
  $\mathcal M_H$ 是超 Kähler 流形：它有三个复结构 $I,J,K$，满足四元数关系。
  \begin{itemize}
    \item 在 $I$ 下：它是稳定 Higgs 丛的模空间；
    \item 在 $J$ 下：它是稳定平坦 $G_\CC$-丛的模空间（Corlette--Donaldson）；
    \item 复维数 $\dim_\CC\mathcal M_H=(2g-2)\dim G$。
  \end{itemize}
  \vspace{0.4em}
  它还是一个完全可积系统：Hitchin 纤维化
  \[
    \mathcal M_H\;\longrightarrow\;\bigoplus_{i=1}^{r}H^0(C,K_C^{\,d_i}),
  \]
  其纤维是 Prym 簇，谱曲线 $\det(\phi-\lambda)=0$ 控制着一切。
  \begin{beizhu}
  \textbf{批注}\quad 第 9 章里 Atiyah--Hitchin 用“谱曲线”描述单极子，用的正是同一个
  Hitchin 系统；只是那里 $C=\RR^2$ 被换成了复曲线，“谱曲线”才成为真正的代数曲线。
  \end{beizhu}
\end{frame}

\begin{frame}{S-对偶 = 镜像对称}
  $\cN=4$ 理论的 S-对偶：
  \[
    \tau\;\longmapsto\;-\frac{1}{n_g\tau},
    \qquad G\;\longmapsto\;\gvee ,
    \qquad \text{电}\;\longleftrightarrow\;\text{磁}.
  \]
  在二维 sigma 模型里，这个对偶变成了
  \[
    \mathcal M_H(G,C)\quad\longleftrightarrow\quad \mathcal M_H(\gvee,C)
  \]
  之间的\textbf{镜像对称}（在 Strominger--Yau--Zaslow 的意义上，纤维是环面）。
  \vspace{0.4em}

  这就是全部秘密：几何 Langlands 对应，是 Hitchin 模空间镜像对称的一个特例。
\end{frame}

\begin{frame}{算子与膜}
  扭化参数（记作 $\Psi$ 或 $t$）决定用 $\mathcal M_H$ 的哪一个复结构，从而决定
  得到的是 A-模型还是 B-模型。
  \vspace{0.4em}
  \begin{itemize}
    \item Wilson 算子与 't Hooft 算子分别实现 $\mathcal M_H$ 上的两类膜；
    \item S-对偶交换电与磁，也就交换这两类膜；
    \item 镜像对称把 A-膜与 B-膜互换。
  \end{itemize}
  \vspace{0.4em}
  于是“Wilson $\leftrightarrow$ 't Hooft”这个物理对偶，变成了
  “A-膜 $\leftrightarrow$ B-膜”这个几何陈述。
\end{frame}

\begin{frame}{特征层是膜，局部系统是点}
  几何 Langlands 的两侧，在 Hitchin 模空间上各是什么？
  \begin{itemize}
    \item \textbf{自守侧}：$\Bun_G$ 上的 Hecke 特征层 $\cF_\cE$。在 $\mathcal M_H$ 上，
      它实现为一个对 Hecke 算子本征的膜（eigenbrane）。
    \item \textbf{Galois 侧}：$C$ 上的 $\gvee$-局部系统 $\cE$。它对应
      $\mathcal M_H(\gvee,C)$ 里的一个点（更准确地说，一个 B-膜）。
  \end{itemize}
  \vspace{0.4em}
  镜像对称把点（B-膜）送到膜（A-膜）。这就是
  “特征层 $\leftrightarrow$ 局部系统”的物理内容。
  \begin{beizhu}
  \textbf{批注}\quad 这一帧是 2008 年 Frenkel--Witten 文章的核心图像之一。
  注意仍然有一环没合上：膜住在 $\mathcal M_H(G,C)$，而 D-模住在 $\Bun_G$；
  两者不是同一个空间，Kapustin--Witten 只在“非分歧”情形、并在某些极限下把两边对上。
  \end{beizhu}
\end{frame}

\begin{frame}{内窥：镜像对称预测数论}
  内窥（endoscopy）原本是 Langlands 函子性里的一类“函数转移”，涉及有相同对偶群的
  两个群。Frenkel 与 Witten（2008）指出：
  \begin{center}
  \begin{tabular}{@{}ll@{}}
    \toprule
    数论侧 & 几何/物理侧 \\
    \midrule
    内窥转移 & 两个 Hitchin 模空间之间的镜像对称 \\
    $L$-包 & 镜像对称下的膜对应 \\
    \bottomrule
  \end{tabular}
  \end{center}
  \vspace{0.4em}
  这不是“用物理解释已知的数学”，而是“用物理\emph{预测}了数学”：镜像对称给出了
  内窥转移的一个几何构造。这是这场演讲之后这条线最有说服力的进展之一。
\end{frame}

\begin{frame}{六维，以及 2024 年}
  \begin{itemize}
    \item Witten（2009）把整个构造提升到六维：六维 $(2,0)$ 理论在
      $\RR\times\Sigma\times C$ 上的边界条件，直接给出几何 Langlands 的两侧，
      不需要再额外“扭化”。
    \item 2024 年，Gaitsgory 与十余位合作者发表五篇论文，证明了非分歧情形下的
      几何 Langlands 猜想（对任意约化群）。
  \end{itemize}
  \vspace{0.4em}
  物理的“证明”仍是启发式的；但数学的证明已经落地，而这个证明用的中心对象
  （仿射 Grassmannian 上的几何 Satake 等价、$\Bun_G$ 上的 Hecke 特征层）正是
  演讲第 6 章介绍的那些。
\end{frame}

\begin{frame}{一张总表}
  \begin{center}
  \renewcommand{\arraystretch}{1.2}
  \begin{tabular}{@{}lll@{}}
    \toprule
    物理 & 群论 & 几何 \\
    \midrule
    电电荷 & $G$ 的表示 & $X^*(T)$ \\
    磁荷 & $\gvee$ 的表示 & $X_*(T)$ \\
    Wilson 算子 & 最高权 & $X^*(T)/W$ \\
    't Hooft 算子 & 余权 & $X_*(T)/W$ \\
    Hecke 修改 & $\gvee$ 的不可约表示 & 支配余权 \\
    S-对偶 & $G\leftrightarrow\gvee$ & $\mathcal M_H(G)\leftrightarrow\mathcal M_H(\gvee)$ \\
    \bottomrule
  \end{tabular}
  \end{center}
  \vspace{0.3em}
  \begin{center}\small
  这张表就是这场演讲的全部内容：四行物理，四行群论，四行几何，说的是同一件事。
  \end{center}
\end{frame}

\begin{frame}{时间线}
  \begin{center}
  \renewcommand{\arraystretch}{1.15}
  \small
  \begin{tabular}{@{}rl@{}}
    \toprule
    1977 & GNO：磁荷是权；Montonen--Olive：电—磁对偶 \\
    1978 & Witten--Olive：中心荷与 BPS 界 \\
    1979 & Osborn：只有 $\cN=4$ 满足判据 \\
    1990 & Beilinson--Drinfeld：几何 Langlands \\
    1994 & Sen 的束缚态；Seiberg--Witten 解 \\
    2005 & Kapustin：'t Hooft 算子的精确化 \\
    2006 & Kapustin--Witten：规范理论解释几何 Langlands \\
    2008 & Frenkel--Witten：几何内窥与镜像对称 \\
    2009 & Witten：六维 $(2,0)$ 构造 \\
    2021 & 本演讲 \\
    2024 & Gaitsgory 等：几何 Langlands 获证 \\
    \bottomrule
  \end{tabular}
  \end{center}
\end{frame}

\begin{frame}{术语速查}
  \begin{center}
  \renewcommand{\arraystretch}{1.15}
  \small
  \begin{tabular}{@{}lll@{}}
    \toprule
    英文 & 中文 & 一句话 \\
    \midrule
    dual group & 对偶群 & 交换特征标格与余特征标格 \\
    magnetic monopole & 单极子 & 磁荷的载体 \\
    't Hooft operator & 't Hooft 算子 & 沿曲线的指定奇点 \\
    Wilson operator & Wilson 算子 & 沿曲线的全纯的迹 \\
    Hecke modification & Hecke 修改 & 只在一点不同的两个丛 \\
    Hecke eigensheaf & Hecke 特征层 & 被 Hecke 算子本征化的层 \\
    geometric Satake & 几何 Satake 等价 & 反常层 $\leftrightarrow$ 对偶群表示 \\
    Hitchin moduli & Hitchin 模空间 & 超 Kähler，可积系统 \\
    BPS bound & BPS 界 & $M\ge|Z|$ \\
    \bottomrule
  \end{tabular}
  \end{center}
\end{frame}

\begin{frame}{延伸阅读}
  \footnotesize
  \begin{itemize}
    \item E. Witten, \emph{Duality, the Langlands Correspondence, and the Influence of
      Michael Atiyah}, INI, 2021（本演讲）。
    \item A. Kapustin, E. Witten, \emph{Electric-Magnetic Duality And The Geometric
      Langlands Program}, Commun. Number Theory Phys. 1 (2007) 1--236, hep-th/0604151.
    \item E. Frenkel, E. Witten, \emph{Geometric Endoscopy and Mirror Symmetry},
      Commun. Number Theory Phys. 2 (2008) 113--283, arXiv:0710.5939.
    \item E. Witten, \emph{Geometric Langlands From Six Dimensions}, arXiv:0905.2720.
    \item D. Gaitsgory et al., \emph{Proof of the geometric Langlands conjecture},
      five papers, 2024.
    \end{itemize}
  \vspace{0.2em}
  \begin{center}\scriptsize
  中译注本与资料卡：数译仓库 \texttt{papers/WIT-2021-ATIYAH-LANG}、
  \texttt{08\_source\_code/SY01-07}（几何内窥）、\texttt{datasets/kw06}。
  \end{center}
\end{frame}

\begin{frame}{}
  \vspace{2.2cm}
  \begin{center}
  {\Large\bfseries 谢谢}\\[0.8em]
  {\small 译文与批注：数译 · 超对称梦想家}\\[0.3em]
  {\small 底本：Edward Witten, INI, 2021-09-23}
  \end{center}
\end{frame}
